% Generated draft: requires mathematical and source review.
\documentclass[11pt]{article}
\usepackage[margin=25mm]{geometry}
\usepackage{fontspec}
\setmainfont{Latin Modern Roman}
\usepackage{amsmath,amssymb,mathtools,unicode-math}
\setmathfont{Latin Modern Math}
\usepackage{enumitem,needspace,xurl,hyperref}
\setlength{\parindent}{0pt}
\setlength{\parskip}{5pt}
\setlength{\emergencystretch}{3em}
\newcommand{\N}{\mathbb{N}}
\newcommand{\Z}{\mathbb{Z}}
\newcommand{\Q}{\mathbb{Q}}
\newcommand{\R}{\mathbb{R}}
\newcommand{\C}{\mathbb{C}}
\newcommand{\F}{\mathbb{F}}
\newcommand{\A}{\mathbb{A}}
\newcommand{\PP}{\mathbb{P}}
\newcommand{\E}{\mathbb{E}}
\newcommand{\eps}{\varepsilon}
\newcommand{\dd}{\,\mathrm{d}}
\begin{document}
\section*{UnsolvedMath ID: 30004788}
\textbf{Source problem code:} catalogue entry
\subsection{Definitions and mathematical statement}
The supplied record asks for a description of a restriction of representations, but does not specify enough hypotheses to determine a unique mathematical formulation. The following gives a conditional conventional interpretation and identifies the remaining qualifications.

Suppose that \(F\) is a nonarchimedean local field, that \(n\geq 1\) is an integer, and that representations are smooth complex representations. These are contextual assumptions for this interpretation; the record does not explicitly supply them. Put \(G_m=\operatorname{GL}_m(F)\), and use the embedding
\[
\iota:G_n\longrightarrow G_{n+1},\qquad g\longmapsto\operatorname{diag}(g,1).
\]
The embedding is also an interpretation requiring confirmation. Let \(\mathcal O_F\) be the ring of integers of \(F\), let \(K_m=\operatorname{GL}_m(\mathcal O_F)\), and let \(I_m\subseteq K_m\) be the inverse image of the upper triangular subgroup under reduction to the residue field.

Under the spherical convention, an unramified representation \(\pi\) of \(G_{n+1}\) satisfies \(\pi^{K_{n+1}}\neq 0\). A generic representation admits a nonzero Whittaker functional: for a nondegenerate smooth character \(\psi\) of the upper unitriangular subgroup \(U_{n+1}\), there is a nonzero linear functional \(\lambda\) such that
\[
\lambda(\pi(u)v)=\psi(u)\lambda(v)
\qquad(u\in U_{n+1},\ v\in\pi).
\]
Here nondegenerate means that the restriction of \(\psi\) to every simple root subgroup is nontrivial. The record does not say whether unramified means spherical or merely having nonzero Iwahori fixed vectors, nor whether \(\pi\) must be irreducible or admissible.

In the smooth complex setting above, let \(\mathcal B_n\) denote the Bernstein component with inertial support given by the diagonal torus of \(G_n\) with its trivial character. This is the usual Iwahori component. Write \(\operatorname{pr}_{\mathcal B_n}\) for projection to that component, and define the candidate object
\[
R_{I_n}(\pi)=\operatorname{pr}_{\mathcal B_n}\bigl(\operatorname{Res}_{\iota(G_n)}^{G_{n+1}}\pi\bigr),
\]
where the restriction is viewed as a \(G_n\)-representation through \(\iota\). The source must be checked to confirm this meaning of restriction to the Iwahori block.

The supplied problem is to determine the isomorphism class of \(R_{I_n}(\pi)\) and construct an explicit model for it. Its independence question asks whether, for every pair \(\pi,\pi'\) in the intended class of unramified generic representations of \(G_{n+1}\), one has
\[
R_{I_n}(\pi)\cong R_{I_n}(\pi').
\]
The intended class and any further restrictions on the pair are unresolved. This is a question, not an asserted isomorphism.

The record additionally asks how products of Steinberg representations behave. Here \(\operatorname{St}_m\) would denote the Steinberg representation of \(G_m\). No product operation, ranks, twists, normalization, or precise restriction assertion is supplied, so this part cannot be stated as a definite quantified claim. The background also refers to a tempered-representation restriction conjecture without giving its formulation or its relationship to the unramified generic question. A tempered hypothesis cannot therefore be inserted into the displayed question as though it were supplied.
\subsection{Short English statement}
Describe the Iwahori component of the restriction from \(\operatorname{GL}_{n+1}(F)\) to \(\operatorname{GL}_n(F)\) for the intended unramified generic representations, and determine whether its isomorphism class is independent of the representation. The source also requests an analysis involving Steinberg products. The field, representation category, meaning of unramified, Steinberg product operation, and relationship to the background tempered conjecture require confirmation.
\subsection{Sources}
\begin{itemize}
\item Supplied local UnsolvedMath registry record, permanent ID \(30004788\), source problem code \texttt{catalogue entry}; registry locations \texttt{data/unsolvedmath/problems.json} and \texttt{catalog/problems.json}. This draft uses the supplied complete record; those files were not independently inspected.
\item Supplied report citation: \emph{Automorphic Forms, Geometry and Arithmetic} (2022), Oberwolfach Reports, DOI 10.4171/owr/2021/39, pp.\ 2089--2156. \url{https://doi.org/10.4171/owr/2021/39}. The record describes Oberwolfach Reports as documenting mathematical workshops and their open-problem discussions. Its proposed year is 2021; no individual proposer is supplied.
\item Supplied literature source, labeled ``original Oberwolfach report'': \url{https://doi.org/10.4171/owr/2021/39}.
\item Supplied literature source, labeled ``Chan, homological branching law'': \url{https://arxiv.org/abs/1905.01668}.
\item Supplied literature source, labeled ``Chan--Savin, branching rules and Hecke algebras'': \url{https://arxiv.org/abs/1605.05130}.
\item Literature assessment dated exactly \texttt{2026-08-22}: the supplied curation labels the problem partially solved and states that Chan and Savin proved important homological branching and projectivity results, while the full tempered-representation restriction conjecture for the Iwahori component remains unproved in general. This is supplied dated curation, not an independently verified assessment. The record expressly requests independent verification of time-sensitive status and cited claims.
\item The supplied statement status is \texttt{corrected\_verified}; its verification note concerns dollar delimiters and preservation of supplied mathematical content. This label does not establish mathematical correctness or resolve the formulation issues above.
\end{itemize}
\end{document}
